\section{Unconditional two-active AA contraction and the cap-SCC
obstruction}

\textbf{Proof-first candidate, 2026-08-12 PDT.} This note separates two
claims which must not be conflated.

\begin{enumerate}
\def\labelenumi{\arabic{enumi}.}
\tightlist
\item
  The proposed finite graph on cap descriptors is not a valid Markov
  reward quotient. A cap records source availability, not the exact
  bounded population, and a directed support graph carries neither
  transition probabilities nor reward occupation weights.
\item
  No such graph is needed for the two-active part of the exact
  18,496-pair remainder. Every affine-feasible two-active descriptor is
  available/available (AA), and the available-target Bellman episode can
  be run \textbf{without stopping at a cap or tier change}. It is
  therefore an unconditional common-marked-potential episode.
\end{enumerate}

The finite identity used in the second claim is frozen at

\begin{CodeBlock}
src/remaining_18496_all_feasible_two_active_aa_certificate.py
SHA-256 25e5c2fce812d2e3d3f02c0c9377533cf16c68313b77bc0d1e7a485052bd68ee

tests/test_remaining_18496_all_feasible_two_active_aa_certificate.py
SHA-256 2445edcbcf66ff2204c821baec455b4f8a416180360a39a6d757712af5fc22e0
\end{CodeBlock}

It enumerates support/descriptor incidences only. It does not enumerate
orientations, rate vectors, populations, paths, or stochastic histories.

\subsection{1. Why a directed SCC is not yet a Foster
corrector}

Consider first a genuine finite Markov-additive kernel. Let \(I\) be a
finite phase set, \(P=(P_{ij})\) a stochastic matrix, and \(r_i\) the
expected reward of one episode begun in phase \(i\). A bounded phase
corrector \(a:I\to\mathbb R\) with

\[
 r_i+\sum_jP_{ij}a_j-a_i\le-\delta                 \tag{1.1}
\]

can exist only if every stationary law \(\pi\) of every closed recurrent
class satisfies

\[
                         \sum_i\pi_i r_i\le-\delta . \tag{1.2}
\]

This is immediate after multiplying (1.1) by \(\pi_i\) and summing. The
converse, with a possibly smaller \(\delta>0\), is the finite
Poisson/Farkas alternative: solve the centered Poisson equation on each
recurrent class and then extend the corrector backwards over the
transient condensation graph.

Thus the load-bearing finite datum is the \textbf{rate-weighted reward
kernel}, not its directed support. The assertion that a closed SCC
merely contains a coercive node is insufficient unless one also proves
all of the following.

\begin{itemize}
\tightlist
\item
  The quotient phase after an episode is determined by the node and the
  episode endpoint.
\item
  Every edge used to reach the coercive node has a uniformly positive
  transition probability on the relevant escaping family.
\item
  The expected positive reward accumulated before that hit is uniformly
  integrable.
\item
  The negative reward at the coercive hit dominates that complete
  prefix.
\end{itemize}

For a state-dependent family of kernels the exact replacement for (1.2)
is negativity against every invariant occupation measure of every
limiting rate/reward kernel. A graph-only SCC test discards precisely
this measure.

\subsection{2. The cap quotient is not
Markov}

The descriptor cap \(2\) means only ``at least two,'' because binary
sources cannot distinguish \(2\) from a larger population for
enablement. It does not mean that the bounded coordinate equals two.

This failure occurs inside the exact 18,496-pair family. Use the ordered
supports

\[
 L_b=\{0,A,2A,A+B\},\qquad
 L_0=\{C,2C,A+C\},                                  \tag{2.1}
\]

and the one-active descriptor with active species \(B\), weight
\((0,1,0)\), and \(A\)-cap zero. The exact remainder certificate
contains its B/F0 failure rows at each \(C\)-cap \(0,1,2\). Choose any
strong orientation of \(L_0\) containing the label \(2C\to C\), for
example the cycle

\[
                         C\longrightarrow A+C
                          \longrightarrow2C\longrightarrow C.       \tag{2.2}
\]

At populations

\[
                         (A,B,C)=(0,N,k),\qquad k\ge2,                \tag{2.3}
\]

the same descriptor node records \(C\)-cap \(2\). If the label
\(2C\to C\) fires, then

\[
 k=2\quad\Longrightarrow\quad\hbox{cap }1,
 \qquad
 k\ge3\quad\Longrightarrow\quad\hbox{cap }2.         \tag{2.4}
\]

Hence the endpoint cap is not a function of the cap node and the
reaction label. Moreover, within the same cap node the \(2C\)-source
propensity is proportional to \(k(k-1)\), whereas the \(C\)-source
propensity is proportional to \(k\). The all-clock transition
probabilities are not functions of the finite cap node either.

Refining cap \(2\) to exact \(k\) repairs Markovianity only by producing
a countable phase. Sending \(k\to\infty\) is not another bounded-cap
node: it is promotion of \(C\) to the active set and requires a new tier
compactification. There is no fixed population threshold at which this
sequence-defined promotion occurs.

The kinetic obstruction is real rather than cosmetic. An accelerated
immigration--death cofactor can leave every fixed box before a slow
service clock with probability tending to one, even though
countable-phase averaging gives a uniform physical-time service
probability. Thus finite box exits may form an apparent exit-only
circulation while the full countable phase is recurrent. A cap-SCC
calculation cannot replace the needed killed countable-phase resolvent
or Poisson corrector.

We therefore record:

\begin{quote}
\textbf{Proposition 2.1 (finite cap-SCC proposal: fail).} A node
recording only the active mask, tier partition, caps, and mark does not
define a finite Markov reward quotient. Consequently a certificate that
every closed directed SCC contains a locally coercive node would not, by
itself, prove a bounded Foster corrector or an unconditional macro.
\end{quote}

This proposition does not say that the networks are transient. It says
that the proposed finite object forgets the rate-weighted occupation
needed to prove recurrence.

\subsection{3. The exact all-feasible-two-active AA
identity}

For the exact outside-mixed remainder, the frozen finite certificate
gives

\[
\begin{array}{c|r}
\text{quantity}&\text{count}\\ \hline
\text{remainder pairs}&18{,}496\\
\text{pairs with a feasible two-active descriptor}&18{,}322\\
\text{affine-feasible two-active incidences}&1{,}140{,}984\\
\text{corrected-cut passes}&1{,}137{,}900\\
\text{corrected-cut failures}&3{,}084\\
\text{incidences with an S-classified linkage}&0.
\end{array}                                                     \tag{3.1}
\]

The unordered available-kind histogram is

\[
\begin{array}{c|rrrrrr}
\text{kind}&Q/Q&Q/U&C/Q&C/U&U/U&C/C\\ \hline
\text{count}&639{,}468&321{,}858&165{,}870&9{,}360&4{,}212&216.
\end{array}                                                     \tag{3.2}
\]

The incidence and identity fingerprints are respectively

\begin{CodeBlock}
a4c4aa42dadabe7e73d46a690f29fe1d457bf57c5bece422b8c0a3fafe72eb99
bc9195d09dc8717381486ae114fcd59e22786c729d957fe209ceac432deaaac8
\end{CodeBlock}

The analytic use of (3.1) is only that both linkage supports satisfy the
symbolic Q/U/C available-target construction. The counts are not a
stochastic proof.

\subsection{4. Remove the structural-exit test from an AA
path}

Fix a strongly connected orientation and positive labelled rates. Let
\((x_n,t)\) be an escaping marked sequence with active coordinates
\(i,j\) and bounded coordinate \(b\). Thus

\[
                  x_{n,i},x_{n,j}\longrightarrow\infty,
                  \qquad \sup_nx_{n,b}<\infty,
                  \qquad x_n\ge t.                   \tag{4.1}
\]

After a subsequence, fix its complete D-tier type. Suppose the linkage
containing \(t\) is Q-, U-, or C-classified. The symbolic bridge
supplies complexes \(q,c\) in that linkage such that

\[
                  q_b\le c_b,\qquad q\succ_D c.       \tag{4.2}
\]

Here \(q\succ_Dc\) means that the enabled falling-factorial ratio of
\(c\) to \(q\) tends to zero on the fixed tier sequence. Strong
connectivity supplies a simple directed path

\[
              t=y_0\longrightarrow y_1\longrightarrow\cdots
                    \longrightarrow y_m=c.            \tag{4.3}
\]

Run the following episode with every physical clock retained.

\begin{enumerate}
\def\labelenumi{\arabic{enumi}.}
\tightlist
\item
  At \(y_r\), take the next ordinary all-clock jump.
\item
  Continue only if the designated label \(y_r\to y_{r+1}\) fires;
  otherwise stop at the competitor's actual endpoint and actual target.
\item
  On reaching \(c\), take one final ordinary all-clock jump and stop at
  its actual endpoint and target.
\end{enumerate}

There is \textbf{no cap, enabled-support, active-set, or tier
structural-exit test} in this rule.

Indeed, after a successful prefix the exact population is

\[
                           x_{r,n}=x_n-t+y_r.          \tag{4.4}
\]

It dominates \(y_r\), so every designated source is physically enabled
by the preceding actual target. At the pre-final endpoint

\[
                           z_n=x_n-t+c\ge c.           \tag{4.5}
\]

the active coordinates still diverge. Equation (4.2) and (4.5) also give
\(z_{n,b}\ge c_b\ge q_b\), so \(q\) is enabled even if the bounded cap
has changed. Bounded displacement preserves every strict
active-coordinate D-comparison. Therefore

\[
 p_c(z_n)\le {K_c(z_n)_c\over K_q(z_n)_q}
                         \longrightarrow0.            \tag{4.6}
\]

This is the decisive point: a cap change along the successful word does
not remove the comparison source. Intermediate cap and tier labels were
proof bookkeeping, not a physical obstruction.

\subsection{5. Exact all-clock reward}

Carry the actual target after every reaction and put

\[
 F(x,t)=\sum_{r=1}^3\log((x_r-t_r)!),
 \qquad W(x,t)=1+F(x,t).                              \tag{5.1}
\]

For a source \(y\), let

\[
 \lambda_y(x)=K_y(x)_y,\qquad
 \Lambda(x)=\sum_y\lambda_y(x),\qquad
 p_y(x)={\lambda_y(x)\over\Lambda(x)}.               \tag{5.2}
\]

The one-jump marked identity is

\[
 D(x,t):=\mathbb E_{x,t}\Delta F
   =\log p_t-\sum_yp_y\log p_y-\log K_t
                         +\sum_yp_y\log K_y
   \le\log p_t+C_K.                                  \tag{5.3}
\]

Let \(D_{r,n}=D(x_{r,n},y_r)\), let \(a_{r,n}\) be the probability of
the designated label at stage \(r\), and let \(J_{r,n}\) be the expected
remaining stopped reward. The literal all-clock recursion is

\[
                 J_{m,n}=D_{m,n},\qquad
                 J_{r,n}=D_{r,n}+a_{r,n}J_{r+1,n}.    \tag{5.4}
\]

All competing clocks are already in \(D_{r,n}\). At the first prefix
source whose probability tends to zero, (5.3) tends to minus infinity.
Every earlier designated probability stays bounded below after
source-ratio compactification, while the positive tail after the first
rare source is multiplied by \(O(p_{y_r})\to0\). If there is no earlier
rare source, use \(c\) and (4.6). Iterating (5.4) gives

\[
                              J_{0,n}\longrightarrow-\infty.         \tag{5.5}
\]

The episode contains at most ten jumps. Source entropy gives, for every
fixed \(s<\infty\),

\[
                  \sup_n\mathbb E[((\Delta F)^+)^s]<\infty.         \tag{5.6}
\]

At each stage the carried target is enabled, so total hazard is bounded
below by the minimum positive labelled out-rate. All fixed moments of
the physical duration are therefore uniformly bounded.

\subsection{6. Finite-menu state
selection}

For a fixed support pair and fixed orientation, put in one finite menu
every simple target-following path associated with every two-active tier
type, available linkage, actual target, and admissible terminal \(c\).
At a marked state select the path with least exact expected stopped
reward; break ties deterministically.

If uniform coercivity failed on an AA two-active escaping family, a
violating sequence would have a subsequence with fixed active pair, tier
type, actual target, available terminal, and simple path. The
corresponding menu score tends to minus infinity by (5.5), and the
selected minimum score is no larger. This is a contradiction. Hence,
outside a finite subset of every AA two-active escaping family,

\[
 \mathbb E_{x,t}
   [W(X_\tau,T_\tau)-W(x,t)+\eta\tau]\le-1           \tag{6.1}
\]

for a fixed sufficiently small \(\eta>0\). Every episode contains at
least one actual jump, has an actual integrable endpoint, and carries
the same proper \(W\). Episodes concatenate without an entrance toll or
a structural-exit handoff. Since a bounded episode cannot destroy either
of two diverging active coordinates along the fixed escaping sequence,
the corresponding endpoint sequence remains in the same two-active
asymptotic regime. This is a sequential statement, not pointwise
invariance of a cap node.

\begin{quote}
\textbf{Theorem 6.1 (unconditional two-active AA episode).} Fix any of
the 18,496 outside-mixed remainder pairs, any strongly connected
orientations, any positive labelled rates, and any closed population
class. Along every affine-feasible two-active escaping sequence, the
finite-menu all-clock rule above satisfies the unconditional
common-marked-\(W\) estimate (6.1). This includes all 3,084
corrected-cut failure incidences. No cap graph, descriptor-exit
circulation, weighted seam, or orientation enumeration is used.
\end{quote}

This is a complete \textbf{local two-active contract}. A global
pair-recurrence deduction still requires every complementary all-active
and one-active escaping sequence to be covered by compatible rules. In
particular, this theorem does not claim that the one-active complement
has already been closed.

\subsection{7. Exact residual after the
repair}

The theorem removes the complete two-active part of the proposed SCC
problem. All-active feasible descriptors pass the corrected cut and have
no bounded cap. The remaining cap issue is therefore one-active, and in
particular the B/F0 phase begun at a Flat0 mark. There the inactive
process may be a genuine countable recurrent phase; Proposition 2.1
shows why it cannot be replaced by the three availability caps.

A valid completion must use one of the following stronger objects.

\begin{enumerate}
\def\labelenumi{\arabic{enumi}.}
\tightlist
\item
  an unconditional finite killed completion that appends the newly
  enabled B-service before stopping;
\item
  a countable-phase killed resolvent/Poisson corrector, with promotion
  to Theorem 6.1 when an inactive coordinate diverges; or
\item
  a population potential giving direct one-active generator descent.
\end{enumerate}

What is no longer needed is a finite cap-SCC theorem for two-active
charts.
